% Source: https://artofproblemsolving.com/community/c3h2231250_nationals_2017_sprint_problem_25
% Stable identity: aops-topic-2231250-q01
Drop the perpendicular from $A$ to line $BC$, with foot at distance $x$ from $B$ (measured towards $C$), and let $h$ be its length.

\begin{asy}
usepackage("amsmath");
size(6cm);
defaultpen(fontsize(10pt));
pair B=(0,0),C=(10,0); real x=(81-169+100)/20; pair A=(x,sqrt(81-x^2)); pair D=(B+C)/2;
draw(A--B--C--cycle,linewidth(0.9)); draw(A--D,red);
dot(D); label("$A$",A,N);label("$B$",B,SW);label("$C$",C,SE);label("$D$",D,S);
label("$9$",(A+B)/2,W);label("$13$",(A+C)/2,NE);label("$5$",(B+D)/2,S);label("$5$",(D+C)/2,S);
\end{asy}

Then $x^2+h^2=81$ and $(10-x)^2+h^2=169$. Subtracting, $100-20x=88$, so $x=\dfrac{3}{5}$ and $h^2=81-\dfrac{9}{25}=\dfrac{2016}{25}$. Since $BD=5$,
\[AD^2=(5-x)^2+h^2=\dfrac{484}{25}+\dfrac{2016}{25}=\dfrac{2500}{25}=100,\]
so $AD=\boxed{10}$.

(Equivalently, by Apollonius' theorem $AB^2+AC^2=2AD^2+2BD^2$: $81+169=2AD^2+50$, so $AD^2=100$.)
